% Official Beamer format: no overlays or paper screenshots.
\documentclass[aspectratio=169,11pt]{beamer}
\usepackage[utf8]{inputenc}
\usepackage[T1]{fontenc}
\usepackage[english]{babel}
\usepackage{amsmath,amssymb}
\usepackage{booktabs}
\usetheme{default}
\usecolortheme{seahorse}
\setbeamertemplate{navigation symbols}{}
\setbeamertemplate{footline}[frame number]
\setbeamerfont{frametitle}{size=\large}
\title{Reservas bancarias, colateral\\y asignaci\'on de capital}
\author{Alejandro Ventura}

\subtitle{Presentaci\'on de tema: Track B, n\'ucleo A}
\date{7 de octubre de 2026, 08:20 (Lima)}
\begin{document}
\begin{frame}[plain]
 \titlepage
 \begin{center}\footnotesize
 \href{https://github.com/alejandroventuraventurameza-tech/ai-project}
 {\texttt{github.com/alejandroventuraventurameza-tech/ai-project}}
 \end{center}
\end{frame}

\begin{frame}{1. La pregunta y el alcance}
\textquestiondown Cu\'ando una expansi\'on de reservas mejora la asignaci\'on de
capital entre empresas heterog\'eneas?
\medskip
\begin{itemize}
 \item Track B: construir el modelo econ\'omico de la tesis.
 \item A, hoy: un banco, dos empresas y oferta de capital fija.
 \item B, presentaci\'on final: instrumentos de precio y cantidad, encaje y
 transiciones de r\'egimen. Din\'amica patrimonial si es tractable.
\end{itemize}
\medskip
El signo depende de qui\'en recibe capital y de qui\'en est\'a restringido.
La manufactura peruana motiva la pregunta. EEA/ENAHO aportan descripci\'on,
sin identificar el efecto causal de pol\'itica monetaria.
\end{frame}

\begin{frame}{2. El antecedente y el modelo que falta}
\textbf{Gonz\'alez, Nu\~no y Thaler (enero de 2026)} ya estudian
pol\'itica monetaria y mala asignaci\'on con heterogeneidad patrimonial.
\[
 q_tk_t\leq\gamma q_ta_t,\qquad
 \max_{k_t,l_t}\{m_tf_t-w_tl_t-R_tk_t\}.
\]
{\small Versi\'on elegida: ec.~4, p.~7; ec.~6, p.~8. Proposici\'on 1, p.~19;
proposici\'on 2, pp.~20--21. Pruebas B.7.1--B.7.2.
$m_t$: precio real del insumo; $R_t$: costo del capital en el paper.}
\medskip

\textbf{Bianchi y Bigio (2022)}: reservas, dep\'ositos, cr\'edito y
gesti\'on de liquidez interbancaria.
\medskip

\textbf{Nuestro objeto}: dos empresas con rendimientos decrecientes, costo
bancario simplificado y colateral de pago heterog\'eneo.
El costo cuadr\'atico que sigue es propio.
\end{frame}

\begin{frame}{2. Banco: reservas escasas y costo marginal}
\[
 \max_{B,D\geq0}\ tB+r_m m-\rho D-C,
 \qquad B+m=D+E,
\]
\[
 C=cB+\frac{\kappa}{2}(\delta D-m)_+^2,
 \qquad c\geq0,\;\kappa>0,\;0<\delta<1.
\]
\textbf{Liquidez escasa e interior bancario:}
\[
 t=\rho+c+\kappa\delta[\delta(B+m-E)-m].
\]
$B$: cr\'edito. $D$: dep\'ositos. $m$: reservas. $E$: patrimonio bancario.
$t,\rho,r_m$: retornos brutos de cr\'edito, ahorro y reservas.
\medskip

Experimento A: m\'as $m$ con $\rho$ fijo. El hogar pasivo financia dep\'ositos
y ofrece capital. El mercado OTC y el bienestar quedan para una extensi\'on.
\end{frame}

\begin{frame}{2. Empresa: fondos propios, deuda y colateral}
Al inicio se alquila capital; al final se produce y se paga la deuda.
\[
 \max_{k_i,x_i,b_i\geq0}\ A_i k_i^\alpha-tb_i+\rho(n_i-x_i)
\]
\[
 \underbrace{Rk_i=x_i+b_i}_{\text{financiamiento}},\qquad
 0\leq x_i\leq n_i,\qquad
 \underbrace{tb_i\leq h_i}_{\text{pago pignorable}}.
\]
\begin{itemize}
 \item $A_H>A_L>0$, $0<\alpha<1$, $n_i,h_i,K>0$, $R>0$, $t>\rho>0$.
 \item $n_i$: fondos iniciales. $h_i$: activos terminales pignorables
 ex\'ogenos. $b_i$: principal, $tb_i$: pago bruto.
 \item Capital fijo: $k_H+k_L=K$. El precio $R$ ajusta.
\end{itemize}
\medskip
Se usa primero patrimonio: $x_i=\min\{Rk_i,n_i\}$,
$b_i=(Rk_i-n_i)_+$.
\end{frame}

\begin{frame}{3. Condiciones marginales y reg\'imenes}
Sea $M_i=\alpha A_i k_i^{\alpha-1}$ y
$\bar k_i=(n_i+h_i/t)/R$.
\medskip

\begin{tabular}{@{}ll@{}}
\toprule
Configuraci\'on & Condici\'on marginal \\
\midrule
$Rk_i<n_i$: autofinancia & $M_i=\rho R$ \\
$Rk_i=n_i$: quiebre & $\rho R\leq M_i\leq tR$ \\
$n_i<Rk_i<n_i+h_i/t$: deuda interior & $M_i=tR$ \\
$k_i=\bar k_i$: colateral activo & $M_i=tR(1+\mu_i)$, $\mu_i\geq0$ \\
\bottomrule
\end{tabular}
\medskip

La concavidad estricta da una demanda \`unica de capital.
Las comparaciones siguientes exigen desigualdades estrictas y permanencia
del r\'egimen. Los quiebres requieren un an\'alisis por separado.
\end{frame}

\begin{frame}{3. P1: H limitada y L se autofinancia}
\textbf{Condiciones}: $b_H=h_H/t$, $M_H>tR$, $Rk_L<n_L$,
$M_L=\rho R$; primitivas y $K$ fijos.
\[
 \frac{n_H+h_H/t}{R}+\left(\frac{\alpha A_L}{\rho R}\right)^s=K,
 \qquad s=\frac1{1-\alpha}.
\]
\[
 \frac{dk_H}{dt}=-\frac{sk_Lh_H}{t^2R(k_H+sk_L)}<0.
\]
Una ca\'ida local de $t$ lleva capital de L a H.
Como $M_H>tR>\rho R=M_L$, H est\'a subcapitalizada.
\[
 \frac{dY}{dt}=(M_H-M_L)\frac{dk_H}{dt}<0.
\]
\textbf{Resultado candidato}: producto sube y mala asignaci\'on cae,
mientras se conserve este r\'egimen.
\end{frame}

\begin{frame}{3. P2: ambas limitadas, signo condicional}
\textbf{Condiciones}: $b_i=h_i/t$, $M_i>tR$ para ambas,
primitivas y $K$ fijos.
\[
 k_H=K\frac{tn_H+h_H}{t(n_H+n_L)+h_H+h_L},
\]
\[
 \frac{dk_H}{dt}=\frac{K(n_Hh_L-n_Lh_H)}{[t(n_H+n_L)+h_H+h_L]^2}.
\]
Al caer $t$, H recibe m\'as capital si y solo si
\[
 \frac{h_H}{n_H}>\frac{h_L}{n_L}.
\]
\textbf{Eficiencia}: requiere adem\'as $M_H>M_L$, equivalente a $k_H<k_H^*$,
y no sobrepasar $k_H^*$ en una intervenci\'on finita.
La igualdad de ratios da neutralidad; el orden inverso revierte la reasignaci\'on.
\end{frame}

\begin{frame}{3. Reservas, equilibrio bancario y eficiencia}
En P1: $B=h_H/t$. En P2: $B=(h_H+h_L)/t$. En ambos, $B'<0$.
\[
 \frac{dt}{dm}=\frac{-\kappa\delta(1-\delta)}{1-\kappa\delta^2B'(t)}<0
 \quad\text{si escasea liquidez}.
\]
Con reservas abundantes, $t=\rho+c$: m\'as $m$ es neutral en este cierre.
\[
 k_H^*=K\frac{A_H^s}{A_H^s+A_L^s},\qquad
 \mathcal L=1-\frac{Y}{Y^*},\qquad G=|\log M_H-\log M_L|.
\]
Una mejora debe elevar $Y$ y acercar los productos marginales.
El ratio de colateral decide la direcci\'on del capital; la brecha inicial
decide si esa direcci\'on mejora eficiencia.
\end{frame}

\begin{frame}{4. B\'usqueda y contribuci\'on por evaluar}
\begin{itemize}
 \item Antecedente principal: enero 2026, 173 p\'aginas PDF. Ec.~4 y 6,
 proposiciones 1--2 y pruebas B.7.1--B.7.2.
 \item Gonz\'alez, Nu\~no y Thaler reconocen a Albrizio y retiran el bloque
 emp\'irico anterior. Versiones 2023 y BdE 2021: contexto hist\'orico.
 \item Bianchi--Bigio: versi\'on local marzo 2021, publicada en 2022.
 Revisi\'on de balances y gesti\'on de liquidez.
 \item Trabajo relacionado localizado: \emph{A tale of two margins}
 (European Economic Review, 2026). B\'usqueda de citantes a\'un parcial.
\end{itemize}
\medskip
\textbf{Aporte por evaluar}: condiciones transparentes para reasignaci\'on
con colateral de pago y dos reg\'imenes. La relaci\'on dinero--mala asignaci\'on
ya es conocida. A no replica su din\'amica ni Ramsey; la novedad sigue en revisi\'on.
\end{frame}

\begin{frame}{5. Plan, Lean y riesgo principal}
\begin{itemize}
 \item \textbf{A}: derivaci\'on completa, KKT en cuatro regiones, equilibrio
 bancario y real; sorteos, casos l\'imite y contraejemplos reproducibles.
 \item \textbf{B, 30 de octubre}: precio, reservas y encaje, transiciones
 financieras; acumulaci\'on patrimonial si el cierre es tractable.
 \item \textbf{Lean despu\'es}: fijar el PDF del art\'iculo en un commit;
 AppliedModelingLib, \texttt{gpt-5.6-sol/xhigh}, carpeta
 \texttt{Ventura26ReservesAllocation}; chequeo \texttt{--fast}.
 \item Copiar corrida completa a \texttt{lean/}; registrar declaraci\'on,
 archivo, commit e hip\'otesis de cada resultado. Sin pruebas Lean terminadas hoy.
\end{itemize}
\medskip
\textbf{Riesgo}: colateral ex\'ogeno y cambios de r\'egimen pueden romper el
signo. La simulaci\'on verifica casos, no demuestra causalidad ni reemplaza la prueba.
\end{frame}
\end{document}
